This review makes a case for legislators who shape medical and
artificial intelligence law, but have little to no robotics or frontier large language model (LLM) application experience. It is built
on a single mechanism, verification before generation, in which a software agent
proposes a clinical action and a ten-gate verification, validation, and
uncertainty quantification examination either accepts it, escalates it to a
qualified human, or blocks it before it can reach a patient. The argument is
organized as eight emotional pillars that legislative-advocacy research finds most
persuasive: compassion, fear of preventable harm, moral outrage, hope,
responsibility, protection of vulnerable people, trust, and urgency. Each pillar
pairs a human appeal with a credible, cited fact. The review draws on a documented engineering lineage,
including a surgical-humanoid assurance run that passed 172 of 172 automated tests
across a ten-gate suite, and on the published advocacy literature describing how
testimony, coalition building, and policy entrepreneurship move a bill through
markup and reconciliation. The conclusion is that Physical
AI Trial Bill H. R. 9510 2026 should be enacted into Federal law.
